\section{Tips}

\subsection{Structure}

A well-structured thesis is essential for effectively presenting your research and arguments. The following guide outlines the typical structure of a thesis, though you should always consult your institution's specific guidelines and requirements. Bellow follows an example structure outline.

\begin{enumerate}

    \item Title Page: The title page should include the title of your thesis, your name, your advisor's name, the name of your institution, and the date of submission. The title page in this template is currently a placeholder, please use the word-template frontpage in your final submission. 

    \item Abstract: The abstract is a concise summary of your research, usually between 150-300 words. It should provide an overview of your research question, methodology, key findings, and conclusions.

    \item Table of Contents: The table of contents lists all the sections and chapters in your thesis, along with their corresponding page numbers.

    \item List of Figures and Tables: If applicable, include a list of all figures and tables in your thesis, along with their corresponding page numbers.

    \item Introduction: The introduction sets the stage for your research, providing the context and background information necessary for understanding your work. It should introduce your research question, explain its significance, and outline your objectives and hypotheses.

    \item Literature Review: The literature review is a critical analysis of the existing research relevant to your topic. It should demonstrate your understanding of the current state of knowledge in your field and identify any gaps or areas of debate that your research will address.

    \item Methodology: In this section, you should describe the research methods and techniques you used to collect and analyze your data. Be sure to provide sufficient detail so that others can replicate your study, and justify your chosen methods in terms of their appropriateness and validity.

    \item Results: The results section presents your research findings, usually in the form of tables, figures, or charts. It should provide a clear and objective account of your data, without any interpretation or discussion.

    \item Discussion: The discussion section is where you interpret and analyze your results, relating them back to your research question and the existing literature. Explain the significance of your findings, discuss any limitations or weaknesses in your study, and suggest areas for future research.

    \item Conclusion: The conclusion provides a summary of your research, restating your main findings and their implications. It should also emphasize the significance of your work, and offer recommendations or suggestions for further research or practice.

    \item References: Include a comprehensive list of all the sources you cited in your thesis, following a consistent citation style (e.g., IEEE, APA, or Harvard).

    \item Appendices: The appendices contain any supplementary material that is not essential to understanding the main text but may be useful for readers seeking additional information. This can include raw data, interview transcripts, or detailed descriptions of research instruments.

\end{enumerate}

\subsection{Work Flow}

I find using the todonotes package is really helpful for my structure and workflow, 
I usually fill in orange todos section-by-section as I get ideas on what needs to be included in a section. While I use yellow boxes to indicate that some text is a draft and still needs to be iterated on. 

examples:
\todo[inline]{for this section, I have to include the following:
\newline I have to include A
\newline I have to include B
\newline I have to include C
}

\paragraph{Text}
\todo[inline, color=yellow]{I have written this text, but only a first draft}
Here is some more \todo{Comment: use the word text or paragraph, make your language consistent} text, I have used some words in this paragraph.

\paragraph{Literature review}

\todo[inline, color=SkyBlue]{Darcy–Weisbach friction factor:
\newline \parencite{Haaland:1983}
\newline\newline Reliability:
\newline \parencite{Rausand:2014} 
\newline \parencite{Rausand:Høyland:2004} 
}

When doing literature reviews as I am reading, I like to add the references I find relevant or interesting like this, which makes it easier to write the review later.

\subsection{Academic Writing}

\begin{enumerate}
    \item Understand the purpose and audience: Before you start writing, make sure you have a clear understanding of the purpose of your paper and the target audience. Consider the expectations and requirements of your readers, and tailor your writing accordingly.

    \item Outline your ideas: Develop an outline to organize your ideas and ensure a coherent structure. This will help you present your arguments in a logical and organized manner, making it easier for readers to follow your train of thought.

    \item Use clear and concise language: Write in a clear and concise manner, avoiding unnecessary jargon and complex sentences. Ensure that your ideas are expressed clearly and effectively, without ambiguity.

    \item Adopt an objective tone: Maintain an objective and unbiased tone throughout your writing. Avoid using emotional language or making unsupported claims. Present your arguments and evidence in a balanced and fair manner, allowing readers to draw their own conclusions.
    
    \item Cite your sources properly: Give proper credit to the sources of your information and ideas by using an appropriate citation style, such as IEEE, APA, or Harvard. This not only ensures academic integrity but also helps readers locate the original sources for further reading.

    \item Use evidence to support your arguments: Back up your claims with relevant and reliable evidence from credible sources. This can include data, research findings, or expert opinions. Make sure to explain the significance of the evidence and how it supports your argument.
    
    \item Be critical: Critically analyze your sources, arguments, and evidence. Identify any potential weaknesses, limitations, or alternative interpretations, and address them in your writing. This demonstrates a deep understanding of the topic and strengthens your argument.

    \item Revise and edit: Always revise and edit your work to ensure it is free of grammatical errors, typos, and inconsistencies. This will help you present a polished and professional piece of writing that reflects the quality of your research and ideas.
    
    \item Seek feedback: Share your work with peers, mentors, or instructors, and ask for their feedback. This can help you identify areas for improvement, as well as provide valuable insights and perspectives that you may not have considered.

    \item Practice: Finally, remember that good academic writing takes practice. Continuously work on honing your writing skills, and be open to learning from your experiences and feedback from others.
    
\end{enumerate}